% =========================================================
% Proof: Self Additive Inverse (Characteristic 2)
% Lemma: Self Additive Inverse — Boolean Rings and Boolean Algebras
% Label: prf:boolean-ring-self-additive-inverse
% =========================================================

\newpage
\phantomsection
\label{prf:boolean-ring-self-additive-inverse}

\begin{remark*}[Return]
\hyperref[lem:self-additive-inverse]{Return to Lemma}
\end{remark*}

\begin{lemma*}[Self Additive Inverse]
In any Boolean ring $R$,
\[
p + p = 0 \qquad \text{for all } p \in R.
\]
\end{lemma*}

\begin{proof}
\textbf{Professional Standard Proof.}~\\
Apply idempotence to the element $p + p$:
\[
(p + p) \cdot (p + p) = p + p. \tag{i}
\]
Expand the left side using right distributivity~(9), then left
distributivity~(8), then collapse each $p \cdot p = p$ by
idempotence~(11):
\[
(p + p) \cdot (p + p)
  \;\overset{(9)}{=}\; p\cdot(p+p) + p\cdot(p+p)
  \;\overset{(8)}{=}\; (p^2 + p^2) + (p^2 + p^2)
  \;\overset{p^2=p}{=}\; (p + p) + (p + p). \tag{ii}
\]
Combining (i) and (ii): $(p+p) + (p+p) = p+p$.
Setting $s = p+p$, the equation $s + s = s$ gives $s = 0$
by adding $-s$ and applying associativity, the inverse law,
and the identity law.
Hence $p + p = 0$.
\end{proof}

\begin{proof}
\textbf{Detailed Learning Proof.}~\\

\textbf{Step 1. Apply idempotence to the element $p + p$.}
The element $p + p$ is a member of the carrier of $R$.
Axiom~(11) (multiplicative idempotence) asserts every element
of $R$ is its own square.
Apply this to the element $p + p$:
\[
(p + p) \cdot (p + p) = p + p. \tag{i}
\]
This is the pivot of the proof.
It says $p + p$ squares to itself.
Nothing about $p + p = 0$ has been assumed;
that is the conclusion to be derived.

\textbf{Step 2. Expand $(p+p)\cdot(p+p)$ using right distributivity.}
Axiom~(9) (right distributivity) states:
\[
(q + r) \cdot t = q \cdot t + r \cdot t
\qquad \forall\, q, r, t \in R.
\]
Set $q = p$, $r = p$, $t = p + p$:
\[
(p + p) \cdot (p + p) = p \cdot (p + p) \;+\; p \cdot (p + p). \tag{ii}
\]
The product has split into two identical copies of $p \cdot (p+p)$.
This step holds in any ring; no idempotence has been used.

\textbf{Step 3. Expand each $p\cdot(p+p)$ using left distributivity.}
Axiom~(8) (left distributivity) states:
\[
t \cdot (q + r) = t \cdot q + t \cdot r
\qquad \forall\, t, q, r \in R.
\]
Apply to $p \cdot (p + p)$ with $t = p$, $q = p$, $r = p$:
\[
p \cdot (p + p) = p \cdot p \;+\; p \cdot p.
\]
Apply to both copies in~(ii):
\[
p \cdot (p+p) + p \cdot (p+p)
= (p \cdot p + p \cdot p) + (p \cdot p + p \cdot p). \tag{iii}
\]
The right side contains four occurrences of $p \cdot p$.
This expansion holds in any ring; idempotence has still not
been used.

\textbf{Step 4. Collapse each $p \cdot p$ to $p$ by idempotence.}
Axiom~(11) gives $p \cdot p = p$.
Substitute into all four occurrences in~(iii):
\[
(p \cdot p + p \cdot p) + (p \cdot p + p \cdot p)
= (p + p) + (p + p). \tag{iv}
\]
This is the second and final use of idempotence.

\textbf{Step 5. Chain (i) through (iv).}
Equations (i), (ii), (iii), (iv) combine as:
\[
p + p
\;=\; (p+p)\cdot(p+p)
\;=\; (p+p)+(p+p).
\]
Hence
\[
(p+p) + (p+p) = p+p. \tag{v}
\]

\textbf{Step 6. Conclude $p + p = 0$ by cancellation.}
Let $s = p + p$.
Equation~(v) states $s + s = s$.
Add the additive inverse $-s$ to the right of both sides,
then apply associativity, the inverse law, and the identity law:
\[
s + s = s
\;\implies\; (s + s) + (-s) = s + (-s)
\;\implies\; s + \underbrace{(s + (-s))}_{=\,0} = 0
\;\implies\; s + 0 = 0
\;\implies\; s = 0.
\]
Since $s = p + p$, the conclusion is $p + p = 0$.
\end{proof}

\begin{remark*}[Proof structure]
The proof is direct.
The strategy applies idempotence twice: once to $p + p$ as a
whole (Step~1) and once to $p$ individually (Step~4), with
distributivity (Steps~2--3) connecting the two applications.
The conclusion follows by a standard additive cancellation
(Step~6) using only the group axioms for $(R,+)$.
No circularity arises: the proof never assumes $1 + 1 = 0$
or any consequence of characteristic~$2$.
Both $p + p = 0$ and $1 + 1 = 0$ are consequences of
idempotence; this proof derives the former directly without
appealing to the latter.
The Lean~4 formalisation is \texttt{AdditiveIdempotence} in
\texttt{LRA/VolumeI/BooleanAlgebra/BooleanAlgebraDefinitions.lean}.
\end{remark*}

\begin{remark*}[Dependencies]~\\
\begin{itemize}
  \item \hyperref[def:ring]{Definition of Ring}
        --- axioms~(1) associativity, (5) identity, (7) inverse,
        (8) left distributivity, (9) right distributivity.
  \item \hyperref[def:boolean-ring]{Definition of Boolean Ring}
        --- axiom~(11) multiplicative idempotence,
        applied in Steps~1 and~4.
  \item \hyperref[def:idempotent-element]{Definition of Idempotent Element}
        --- names the property $p \cdot p = p$.
\end{itemize}
\end{remark*}
